\section{Every minimal CI crossing of six or more factors has all-padding late derivative LC}
This is an unbounded derivative-LC entrance theorem. Its seventeen
scoped kernel reports and independent positive-box reconstruction
are recorded in the accompanying verification manifest.
It is not a theorem of general C, A2 or original one-descent UI.
All LC numerators have positive interval support. Put
$q=\phi-1$, $\psi=-q$.

\paragraph{Statement.}
Fix a binary count $u$. Let $V=(t_1,\ldots,t_r)$ be a sorted
deletion-minimal degree crossing, with $t_i\ge3$, $r\ge6$,
$h=\sum(t_i-1)>u$, and $h-\min(t_i-1)\le u$.
For every padding $\alpha\ge1$, the actual positive derivative
tail of $1^\alpha2^uV$, from physical $W=\sum t_i$, is LC
at every internal centre before its first tail mode.
Equivalently it is LC at every pre-global-mode centre $N\ge W+1$.

The crossings $h=u+1,u+2$ are already proved. Hence write
$h=u+\delta$, $\delta\ge3$. Minimality forces EVERY degree
$d_i=t_i-1\ge\delta$, so every capacity is at least four.
Also $(r-1)\delta\le u$, so $u\ge15$.

\paragraph{There are no long endpoint rows.}
Delete the LARGEST capacity $t$ and call the remaining product $L_0$,
degree $h_0$, length $r_0=r-1$. Then
\[
 h_0=u+\delta-t+1\le u,\quad r_0\le u/\delta\le u/3,
 \qquad t\le u-(r-2)\delta+1\le u-11.              \tag{1}
\]
Thus the old quotient belongs to the registered degree-boundary or
surplus family: positive global ordinary LC and positive primitive
LC after its first internal centre. Its only old exceptional quotient
has $u=3$, outside the present range. There is no new endpoint at
$u-1,u,u+1$, and the old weakened uniform-window primitive theorem
applies after centre $t$.

\paragraph{An unscaled interior budget already at sixteen.}
For $u\ge16$, $Q=Q_u$ the registered degree-$u$ truncated quotient,
$2\le s\le u-3$, and $1\le k\le s$, the ULC factor proof in
\texttt{unscaled\_head\_and\_fast\_source\_proof.tex} gives
\[
 K(s,k)\ge(5/4)T_k.
\]
The scalar payment is
\[
 (5/4)(u+2)^2(u+q)\le2(u+1)(u-1)(u-2).
\]
With $q\le2/3$, its margin times twelve is
$9u^3-118u^2-124u+8$, which at $u=16+x$ is
$9x^3+314x^2+3012x+4680>0$.
For ANY interval-supported nonnegative LC $N_0=1$,
$N_1\le u/3+1$, with NO degree bound, the first-intermediate-row
Cauchy--Binet argument therefore proves
\[
 M_{NQ}(s)\ge(5/4)
       (B_{s+1}+2qB_s+q^2B_{s-1})\quad(2\le s\le u-3).
                                                               \tag{2}
\]
Row zero, including $N_{s+1}\le N_1N_s$, is paid by
$\binom{u-1}{2}\ge(5/4)(4u/3+5/3)$.

\paragraph{Six factors make the natural spectral atom uniformly small.}
For EVERY $t\ge4$ the interval norm is at most
$\rho=q+q^6<27/40$: odd capacities are at least five and even
ones have the smaller norm $q-q^{t+1}<q$. The bounds
$q>309/500$, $q<5/8$, and $q^6=5-8q$ suffice. Hence
\[
 \phi^2\rho^r\le(21/8)(27/40)^6<1/4.              \tag{3}
\]
Put $L=L_0[t]$, $c=\psi^uL(-\phi)$,
$J=((1+z)^uL-c)/(1+qz)$, $a=\phi^2c$.
The sign or size of its EARLY coefficients is not assumed.
Its natural tail starts at rank $\delta-1$. There
$|c|q^{\delta-1}\le\phi\rho^r<1/4$, and at every required
ordinary centre $n\ge\delta$,
$|c|q^{n-1}<1/4$. Thus (2), now with $N=L$, proves all ordinary
head minors at centres $\delta..t-1$. The head values are positive.
The correction to $[t]J_{u,L_0}$ has degree $t-2$, so all
ordinary centres at least $t$ are inherited. Consequently the
natural tail $b_k=J_{\delta-1+k}$ is positive and globally LC.

For its primitive take $B=L_0Q_u$, $X_n=B_n-a\psi^n$.
Its exact head reads are
$d_n=X_n$ for $1\le n<t$, $d_t=X_t-1$,
$d_{t+1}=X_{t+1}-B_1$.
For $n\ge\delta+1$, the perturbation coefficient
$|a|q^{n-1}\le\phi^2\rho^r<1/4$. Thus (2) pays every inner
head minor; at $t-1$, lowering the next read by one improves it.
All involved read ranks are in the positive head, by (1)--(3).
The registered $3u/2$ interval-gap theorem applies with
$r_0\le u/3\le u/2$ and gives
$B_1B_{t-1}-2B_t\ge3u/2$.
Its perturbation error is at most
$(1/4)(B_1+2q)\le u/3+1/6$.
This pays the exact last minor
$M_X(t)+1+B_1X_{t-1}-2X_t$.
Every later primitive centre, from $t+1$, is inherited from the
old weak-window theorem. The primitive of $b$ AFTER deleting
rank zero is therefore globally LC: its first relevant centre
is original rank $\delta+1$. No claim is made about its omitted
first primitive centre.

\paragraph{The actual inverse and its full finite source.}
Let $x_s=\sum_{j\ge1}q^jJ_{s+j}$. The positive LC natural tail
makes $x$ LC at every centre $s\ge\delta-1$ and gives
$x_s/x_{s-1}<\phi$ there. The exact identity is
\[
 x_{\delta-2}=C_{h-2}+1,\qquad
 x_s=C_{u+s}\ (s\ge\delta-1),\quad
 C=L\sum_{k=0}^{u-1}(z+z^2)^k.
\]
These are the ACTUAL infinite-initial-ones signed derivatives from
physical $W$, including its previous unit. For $a_0=\alpha-1$,
the finite actual tail is $(x_s-E_{s-a_0})^+$,
$E=(1+z)^uL$. The true spectral atom in
$E=(1+qz)J+c$ is retained.

Apply the algebraic signed-primitive kernel identity with lag
$a_0$, which remains nonnegative. Discard all primitive indices
below $\delta$. Since $K_{a_0}$ vanishes above $a_0$, this
whole discarded prefix has ZERO contribution in output triples
centred at $s\ge a_0+\delta+1$.
The remaining positive primitive is globally LC; its later
negative part is nonpositive and produces exactly the geometric
constant-cut term. The registered constant-cut proof therefore
pays every such triple before the first natural $J$ mode.

Earlier startup triples have source centre $j=s-a_0\le\delta$.
If $j<0$ the source is absent or only lowers the next read.
Source centres zero and one use the registered all-padding
two startup identities; $E_1=u+r\ge21$ here.
All $2\le j\le\delta$ are paid by the fast-source theorem
whenever
\[
                    13(\delta+1)\le3(u+1).         \tag{4}
\]
Indeed $jE_j\ge(u-j+1)E_{j-1}$ and the same next edge give
both source growth ratios at least $10/3\ge2\phi$.
Every background triple is in the proved LC inverse tail.
Finite/infinite same-word positive-derivative density puts the
finite first tail mode no later than the infinite one; the inverse
mode is no later than the first natural quotient mode. Thus every
required pre-mode triple belongs to one of the paid ranges.

\paragraph{Exactly three exceptions, completely finite-covered.}
Since $u\ge5\delta$, failure of (4) forces
$(u,\delta)=(15,3),(16,3),(20,4)$. Their degree sums and
$r\ge6$ force EXACTLY
\[
             V=4^6,\qquad V=4^5\,5,\qquad V=5^6
\]
respectively. The registered padding clamp gives $17,18,23$
hosts, in total 58. All their pre-first-tail-mode minors have
an independent positive suffix-box reconstruction; the 449 distinct
integer margins in the 770 comparisons are kernel checked.
This also covers the sole $u=15$ case below the analytic threshold.
The theorem follows for EVERY minimal crossing with $r\ge6$.

\paragraph{A constant-five-factor entrance reduction.}
Every CI of degree $h>u$ contains a deletion-minimal crossing or
a previously proved long single-factor seed. If that minimal
crossing has at least six factors, the theorem just proved applies.
The registered arbitrary-insertion signed-tail identity propagates
its late derivative LC through every additional CI capacity.
Hence the universal actual late-D-LC problem now needs only
minimal crossings with AT MOST FIVE CI factors. They may have
NO capacity three (the degree-$u+1,u+2$ crossings are already closed),
and still obey $u+3\le h\le2u$, $h-\min d_i\le u$,
$3\le d_i\le u$, with the old analytic padding clamp.
No bounded-count entrance is asserted to have been checked in full.
This is a derivative-LC reduction; earlier F/G, C, A2 and original
one-descent coefficient unimodality remain separate and OPEN.

\paragraph{Formal scope.}
The unscaled-head and clipped-fast-source kernels are reused from
\texttt{UnscaledHeadFastSource.lean}. Seven additional unbounded
scalar/count/prefix-support kernels are in
\texttt{ManyFactorCICrossing.lean}; one complete integer certificate
checks all 449 bounded-exception margins. The quotient identities,
ULC/Cauchy--Binet/window assembly, actual tail alignment and mode
transport are written, not an end-to-end Lean composition theorem.
